\documentclass[12pt,a4paper]{article}
\usepackage[utf8]{inputenc}
\usepackage[T1]{fontenc}
\usepackage[brazil]{babel}
\usepackage{indentfirst}
\usepackage[margin=2.5cm]{geometry}

\title{Resenha Cr\'itica: Artigo 2 ---\\
On the Criteria to Be Used in Decomposing Systems into Modules}
\author{Gustavo Pessoa Firmino Duarte}
\date{8 de Mar\c{c}o de 2026}

\begin{document}
\maketitle

\section{Introdu\c{c}\~ao}

Publicado em 1972 na \textit{Communications of the ACM}, o artigo de David L. Parnas,
``On the Criteria to Be Used in Decomposing Systems into Modules'', \'e um marco
fundamental da engenharia de software. Em uma \'epoca em que a modulariza\c{c}\~ao
era frequentemente confundida com simples divis\~ao de fluxo de execu\c{c}\~ao,
Parnas prop\~oe uma crit\'ica precisa e construtiva: o crit\'erio correto para
decompor um sistema n\~ao \'e a sequ\^encia de passos computacionais, mas sim o
\textit{ocultamento de informa\c{c}\~ao} (\textit{information hiding}). Cinquenta
anos depois, esse princ\'ipio continua sendo um dos pilares do design de software
moderno, presente em arquiteturas orientadas a objetos, microsservi\c{c}os e APIs.

\section{O Princ\'ipio do Ocultamento de Informa\c{c}\~ao}

Parnas utiliza o exemplo do sistema KWIC (\textit{Key Word in Context}) para
contrastar duas estrat\'egias de modulariza\c{c}\~ao. A primeira decomp\~oe o sistema
seguindo o fluxo algor\'itmico: leitura de entrada, c\'alculo das deslocamentos
circulares, ordena\c{c}\~ao e sa\'ida. Cada m\'odulo corresponde a uma etapa do
processo. A segunda decomp\~oe com base no ocultamento de decis\~oes de projeto:
cada m\'odulo esconde uma decis\~ao que pode mudar --- como a estrutura de dados
usada para armazenar as linhas, o algoritmo de ordena\c{c}\~ao ou o formato de
sa\'ida.

O autor demonstra que a primeira abordagem cria um acoplamento alto: qualquer
mudan\c{c}a na representa\c{c}\~ao interna dos dados exige altera\c{c}\~oes em
m\'ultiplos m\'odulos. J\'a na segunda abordagem, cada decis\~ao de projeto fica
encapsulada, tornando as mudan\c{c}as locais e invis\'iveis aos demais m\'odulos.
Parnas defende que um bom m\'odulo deve ter uma \textit{interface} est\'avel que
esconda seus detalhes internos, permitindo que a implementa\c{c}\~ao evolua sem
propaga\c{c}\~ao de impacto.

Al\'em da manutenibilidade, Parnas aponta outros benef\'icios do ocultamento de
informa\c{c}\~ao: compreensibilidade seletiva (\'e poss\'ivel entender um m\'odulo
sem conhecer os demais), desenvolvimento independente por equipes diferentes e
maior facilidade de testes unit\'arios. Esses benef\'icios antecipam, de forma
not\'avel, os princ\'ipios do desenvolvimento \'agil e da arquitetura orientada a
componentes que dominariam as d\'ecadas seguintes.

\section{Conex\~ao com a Pr\'atica Profissional}

A leitura deste artigo ressoa diretamente com minha experi\^encia no desenvolvimento
de um MVP de e-commerce utilizando Node.js, Prisma ORM, React e SQLite. Ao
escolher o Prisma ORM como camada de acesso a dados, tomei --- sem ainda ter
lido Parnas --- exatamente a decis\~ao que ele recomenda: esconder do restante da
aplica\c{c}\~ao os detalhes do banco de dados relacional subjacente. O Prisma age
como um m\'odulo no sentido parnesiano: os servi\c{c}os de dom\'inio apenas chamam
m\'etodos como \texttt{prisma.product.findMany()}, sem qualquer conhecimento de
SQL, \'indices ou estrutura de tabelas.

Quando, em determinado momento do projeto, precisei migrar parte dos dados para
uma estrutura diferente, apenas o schema do Prisma e as migra\c{c}\~oes precisaram
ser alterados --- as camadas de servi\c{c}o e de apresenta\c{c}\~ao permaneceram
intocadas. Essa experi\^encia concreta validou o argumento central de Parnas:
sistemas decompostos por ocultamento de informa\c{c}\~ao s\~ao substancialmente
mais f\'aceis de modificar.

Tamb\'em reconhe\c{c}o, em retrospecto, viola\c{c}\~oes desse princ\'ipio. Em um
projeto acad\^emico de cinco dias, sob press\~ao de prazo, struct\'utei v\'arios
componentes React que consumiam diretamente objetos do banco de dados, sem nenhuma
camada de abstra\c{c}\~ao. Qualquer altera\c{c}\~ao no modelo de dados exigia
busca e substitui\c{c}\~ao em dezenas de componentes --- exatamente o problema
descrito por Parnas na decomposi\c{c}\~ao baseada em fluxo.

\section{Conclus\~ao}

O artigo de Parnas \'e notavelmente atemporal. Sua proposta central --- decompor
sistemas com base nas decis\~oes de projeto que cada m\'odulo esconde --- continua
sendo o crit\'erio mais s\'olido que conhe\c{c}o para avaliar a qualidade de uma
arquitetura. O princ\'ipio do ocultamento de informa\c{c}\~ao n\~ao \'e apenas
uma regra te\'orica; \'e uma ferramenta pr\'atica que reduz o custo de mudan\c{c}a,
diminui o acoplamento e facilita o trabalho em equipe. Em um cen\'ario de software
cada vez mais distribu\'ido e din\^amico, onde decis\~oes tecnol\'ogicas mudam com
frequ\^encia, a lic\~ao de Parnas torna-se ainda mais valiosa: projete para a
mudan\c{c}a, n\~ao apenas para o funcionamento presente.

\end{document}
